\examyear{2012}
\examera{平成24年}
\examfield{解析}
\examgroup{B}
\examnumber{18}
\examtitle{平成24年 B 第18問}
\examtag{固有値問題}
\examtag{微分作用素}
\examtag{特殊関数}

\(\beta>\frac12\) とする。\(x>0\) の領域で定義された関数に作用する微分作用素
\[
  H=-\frac{d^2}{dx^2}+\left(x^2+\frac{\beta(\beta-1)}{x^2}\right)
\]
の固有値問題
\[
  H\psi(x)=E\psi(x)
  \tag{*}
\]
を考える。

\begin{enumerate}
  \item \(\psi_0(x)=x^\beta e^{-\frac12x^2}\) が固有値問題 \((*)\) の解であることを示し，対応する固有値 \(E\) を求めよ。

  以下ではこの固有値を \(E_0\) で表すことにする。
  \item \(\psi(x)=\psi_0(x)\phi(x)\) と置くと，方程式 \((*)\) は \(\phi(x)\) に関する微分方程式
  \[
    \phi''(x)+2\left(\frac{\beta}{x}-x\right)\phi'(x)+(E-E_0)\phi(x)=0
    \tag{**}
  \]
  に帰着されることを示せ。
  \item \(\phi(x)\) が \(x\) の多項式で，恒等的に \(0\) ではなく，かつ微分方程式 \((**)\) の解でもあるという。
  そのような \(\phi(x)\) を固有値 \(E\) とともにすべて決定せよ。
\end{enumerate}